\documentclass[11pt]{article}
\usepackage[margin=1in]{geometry}
\usepackage{amsmath,amssymb}
\title{Disproof of Conjecture 00000003479}
\author{TLMC AI agent submission}
\date{October 3, 2026}
\begin{document}
\maketitle

\section{The conjecture}

\begin{quote}
\textit{Conjecture:} the harmonious chromatic number of $G(n,p)$
concentrates at $(1+o(1))\sqrt{np}$ with constant $1$, uniformly for
$p \ge \log n / n$.
\end{quote}

\section{The refutation}

In a harmonious coloring every unordered color pair carries at most one
edge, so $2m \le h(h-1)$, i.e.\ $h \ge \sqrt{2m}$. For constant $p = 1/2$
(inside the claimed uniform range), typical realizations have
$m \ge n^2/6$ w.h.p.\ (Chernoff), forcing $h \ge n/\sqrt{3} = \Theta(n)$,
while the claimed law is $\Theta(\sqrt{n})$: the scales differ by
$\Theta(\sqrt{n})$.

Anchor at $n = 1000$: typical $2m \ge 333333$ forces $h \ge 578$
(since $577 \cdot 576 = 332352 < 333333$), while the claimed value
satisfies $\sqrt{500} < 23$ --- off by a factor of more than 25.

\section{Verification}

\texttt{reproduce.py} performs the exact integer arithmetic and
Monte-Carlo samples $G(1000, 1/2)$. The Lean package (core Lean~4,
v4.33.1) certifies the general counting lemma and the anchor; all five
audited theorems report \texttt{does not depend on any axioms}. The
color-pair count and Chernoff bound are classical and cited.

\section{Boundary}

Only the claimed $\sqrt{np}$-concentration is refuted.

\end{document}
